% section: 解を組み合わせる見方

\section{解を組み合わせる見方}\label{節：解を組み合わせる見方}
  基本4型では,一つの\bookterm{sequence}{数列}を初項から求めた.
  次は、\bookterm{initial}{初期条件}を変えたときに生まれる解を、まとめて見る。
  同じ式を満たす\bookterm{sequence}{数列}を、定数倍して足す。
  その操作から、隣接多項間・連立型・虚数解型と、付録の\bookterm{matrix}{行列}がつながる。
  まず具体例で組み合わせ方を確かめ,その後に用語を紹介する.

  \begin{bookpoint}
    \textbf{解を組み合わせる}\par
    同じ式を満たす\bookterm{sequence}{数列}を組み合わせ、\bookterm{initial}{初期条件}を合わせる。
    必要な\bookterm{sequence}{数列}の個数は、自由に選べる初期値の個数とつながる。
  \end{bookpoint}

  \subsection{解に定数を掛けて足す}\label{節：線形結合}
    次の項が二つ前の項と同じになる
    \[
      a_{n+2}=a_n\qquad(n\ge1)
    \]
    を考える.初めの二項を$(1,0)$とすると,\bookterm{sequence}{数列}は$1,0,1,0,\ldots$となる.
    $(0,1)$とすると、$0,1,0,1,\ldots$となる。それぞれを$u_n,v_n$と書く。
    \[
      \begin{array}{c|rrrr}
        n&1&2&3&4\\ \hline
        u_n&1&0&1&0\\
        v_n&0&1&0&1
      \end{array}
    \]
    例えば$3u_n-2v_n$は$3,-2,3,-2,\ldots$であり,これも同じ\bookterm{recurrence}{漸化式}を満たす.
    一般の定数$C,D$についても,
    \[
      Cu_{n+2}+Dv_{n+2}=Cu_n+Dv_n
    \]
    だから,$Cu_n+Dv_n$は解である.
    このように\bookterm{sequence}{数列}に定数を掛けて足したものを,それらの\textbf{\bookterm{combination}{線形結合}}という.

    この性質を使うときは,\textbf{\bookterm{initial}{初期条件}を固定する前の解全体}を考えている.
    組み合わせた\bookterm{sequence}{数列}の初めの二項は$(C,D)$なので,\bookterm{initial}{初期条件}は係数によって変わる.

    同じことは,\bookterm{homogeneous}{同次}\bookterm{linear}{線形}\bookterm{recurrence}{漸化式}に共通して成り立つ.
    例えば$u_n,v_n$がともに
    \[
      w_{n+2}+pw_{n+1}+qw_n=0
    \]
    の解なら,$b_n=Cu_n+Dv_n$について
    \BookMathLayout{\[
      \begin{aligned}
        b_{n+2}+pb_{n+1}+qb_n
        &=C(u_{n+2}+pu_{n+1}+qu_n)\\
        &\quad+D(v_{n+2}+pv_{n+1}+qv_n)\\
        &=0
      \end{aligned}
    \]}{\[
      \begin{aligned}
        b_{n+2}+pb_{n+1}+qb_n\\
        {}=C(u_{n+2}+pu_{n+1}+qu_n)\\
        \quad+D(v_{n+2}+pv_{n+1}+qv_n)\\
        {}=0
      \end{aligned}
    \]}
    である.掛けて足す操作と更新が両立するためである.
    この性質を\textbf{\bookterm{superposition}{重ね合わせの原理}}という.
    \bookterm{homogeneous}{同次}・\bookterm{inhomogeneous}{非同次}の意味は\sectionref*{節：アフィン・同次・非同次}で確認できる.

  \subsection{必要な解の数と,初期条件の自由度}\label{節：一次独立・基底・次元}
    先ほどの$u_n,v_n$を使えば,どんな\bookterm{initial}{初期条件}$(a_1,a_2)$にも
    \[
      a_n=a_1u_n+a_2v_n
    \]
    と合わせられる.右辺は元の\bookterm{recurrence}{漸化式}を満たし,初めの二項も一致する.
    直前の二項から次の項が一意に決まるので,その後のすべての項も一致する.
    また,$Cu_n+Dv_n=0$がすべての$n$で成り立つなら,
    $n=1,2$を代入して$C=D=0$を得る.係数の選び方にも重複がない.

    二つの解を並べるだけで,必ずこのようになるわけではない.
    例えば$a_{n+1}=2a_n$の解$u_n=2^{n-1}$と$v_n=2\cdot2^{n-1}$では,
    \[
      Cu_n+Dv_n=(C+2D)2^{n-1}
    \]
    であり,二つの係数の選び方は一意ではない.一つの係数にまとめられる.

    ここまでの違いを表す名前もある。
    \bookterm{sequence}{数列}の組
    \[
      u^{(1)}_n,\ldots,u^{(m)}_n
    \]
    について,
    \[
      C_1u^{(1)}_n+\cdots+C_mu^{(m)}_n=0
    \]
    がすべての$n$で成り立つのは$C_1=\cdots=C_m=0$の場合だけ,という関係を\textbf{\bookterm{independent}{一次独立}}という.
    二つの非零の\bookterm{sequence}{数列}の場合には,一方が他方の定数倍にならないことに当たる.
    \bookterm{independent}{一次独立}な解の組で,すべての解を\bookterm{combination}{線形結合}として表せるものを\textbf{\bookterm{basis}{基底}}という.
    \bookterm{basis}{基底}を選ぶと,各解を表す係数は一意に決まる.
    その\bookterm{basis}{基底}に必要な\bookterm{sequence}{数列}の個数を\textbf{\bookterm{dimension}{次元}}という.
    $a_{n+2}=a_n$の解全体は2\bookterm{dimension}{次元}であり,$a_{n+1}=2a_n$の解全体は1\bookterm{dimension}{次元}である.

    \bookterm{homogeneous}{同次}\bookterm{linear}{線形}\bookterm{recurrence}{漸化式}の解全体には零\bookterm{sequence}{数列}が含まれ,足し算や定数倍をしても解のままである.
    こうした操作のできる集合を\textbf{\bookterm{vector-space}{ベクトル空間}}という.
    \bookterm{sequence}{数列}も,矢印と同じように足したり定数倍したりする対象として扱えるのである.
    ここでは係数を実数としている.後の虚数解型で複素数を使うときも,
    係数を複素数にして同じ計算と判定を使える.どちらの係数で考えるかは区別する.

    この「\bookterm{dimension}{次元}」は,更新ができる範囲で\bookterm{initial}{初期条件}を自由に選ぶ数とつながっている.
    例えば,すべての$n\ge1$で係数が定義された\bookterm{homogeneous}{同次}\bookterm{linear}{線形}\bookterm{recurrence}{漸化式}
    \[
      a_{n+k}=\sum_{j=0}^{k-1}c_j(n)a_{n+j}
    \]
    では,初めの$k$項から次の項が一意に決まる.
    初めの$k$項のうち一つだけを$1$,残りを$0$とした$k$個の解を作れば,
    それらの\bookterm{combination}{線形結合}でどんな\bookterm{initial}{初期条件}にも合わせられる.
    係数の一意性も初めの$k$項を比べれば分かるので,解全体は$k$\bookterm{dimension}{次元}である.
    \bookterm{initial}{初期条件}を指定した後は,その中から一つの解を選んでいる.
    この結論は、\bookterm{homogeneous}{同次}\bookterm{linear}{線形}の式で、更新が続けられることを使っている。
    \bookterm{nonlinear}{非線形}の式や、途中で更新できなくなる式では、その条件から調べる必要がある。

  \subsection{同じ更新をする量を引くと,何が残るか}\label{節：非同次と解の組合せ}
    $a_{n+1}=2a_n+3$の二つの解を足すと,次の項には付加項$6$が現れる.
    従って,\bookterm{inhomogeneous}{非同次}の式では任意の\bookterm{combination}{線形結合}がそのまま解になるわけではない.
    一方,同じ更新をする二つの解$a_n,g_n$を引けば,
    \[
      (a_{n+1}-g_{n+1})=2(a_n-g_n)
    \]
    となる.共通の付加項が消え,\bookterm{homogeneous}{同次}の式が残る.
    この例では$g_n=-3$を選べるから,
    \[
      a_n=-3+C2^{n-1}
    \]
    と,一つの既知の解に\bookterm{homogeneous}{同次}の解を加えて全体を表せる.
    \sectionref*{節：補助数列の設計}で引いた「同じ更新をする量」は,
    解全体をこのように分ける役割も果たしていた.

    この見方は,後の節で次のように使う.
    \begin{itemize}
      \item \sectionref*{節：階差等比複合型}では,同じ付加項をもつ解を一つ見つけ,引き算で\bookterm{homogeneous}{同次}に直す.
      \item \sectionref*{節：隣接3項間漸化式型}では,\bookterm{root}{特性根}から独立な二つの解を作り,\bookterm{initial}{初期条件}を合わせる.
            重解のときは同じ冪を二つ並べても足りないので,別の形の解を探す.
      \item \sectionref*{節：連立型}では,二つの量の和と差を使って更新を分ける.
            \sectionref*{節：虚数解型隣接3項間漸化式}では,複素数係数で使う冪の\bookterm{basis}{基底}を,
            実数係数で使える正弦・余弦の\bookterm{basis}{基底}へ直す.
    \end{itemize}
    ある解をほかの解から組み立てることと,複数の項を合わせて\bookterm{auxiliary}{補助数列}を作ることは,
    どちらも「定数を掛けて足す」操作を使う.前者は解全体を表すため,
    後者は更新を簡単にする量を取り出すために使っている.
    付録の\sectionref*{節：行列の基本}では,この二つの役割をベクトルと\bookterm{matrix}{行列}から見直す.
